%\part*{Euclidean geometry}
\addtocontents{toc}{\protect\contentsline{part}{\protect\numberline{}Geometri Euklides}{}{}}
  
\chapter{Aksioma}
\label{chap:axioms}

\vspace*{-3\baselineskip}

Suatu sistem aksioma sudah terdapat dalam ``Elements'' karya Euklides --- buku pelajaran paling berhasil dan berpengaruh yang pernah ditulis.

Kajian sistematis atas berbagai geometri sebagai sistem aksiomatik
 dipicu
oleh penemuan geometri non-Euklides.
Cabang matematika yang muncul dengan cara ini disebut landasan geometri.

Sistem aksioma yang paling populer
diusulkan oleh David Hilbert pada tahun 1899.
Sistem ini juga merupakan sistem pertama yang ketat menurut standar modern.
Sistem tersebut terdiri atas dua puluh aksioma dalam lima kelompok, enam ``pengertian pangkal'', dan tiga ``istilah pangkal'';
semuanya tidak didefinisikan berdasarkan konsep-konsep yang telah didefinisikan sebelumnya.

Kemudian, banyak sistem lain diusulkan.
Beberapa yang patut disebut ialah
sistem yang dikembangkan oleh Alexandr Alexandrov \cite{alexandrov}, yang intuitif dan elementer,
sistem karya Friedrich Bachmann \cite{bachmann} yang didasarkan pada konsep simetri,
serta sistem karya Alfred Tarski \cite{tarski} --- suatu sistem minimalis yang dirancang untuk dianalisis menggunakan logika matematika.

Kita akan menggunakan sistem lain yang dekat dengan sistem George Birkhoff \cite{birkhoff}.
Sistem ini bertumpu pada pengamatan utama (\ref{preaxiomI})--(\ref{preaxiomV}) yang tercantum pada Bagian~\ref{preaxioms}.

Aksioma-aksioma itu memakai gagasan 
ruang metrik, 
garis, 
sudut,
segitiga,
kesamaan modulo $2\cdot\pi$ ($\equiv$), 
kekontinuan pemetaan antara ruang-ruang metrik,
dan kekongruenan segitiga ($\cong$).
Semua konsep ini dibahas dalam bagian pendahuluan.

Sistem kita dibangun di atas ruang metrik.
Secara khusus, kita menggunakan bilangan real sebagai unsur pembangun. 
Karena itu, pendekatan kita tidak sepenuhnya aksiomatik --- kita membangun teori ini di atas sesuatu yang lain;
pendekatan ini menyerupai pengantar geometri Euklides berbasis model yang dibahas pada Bagian~\ref{page:model}.
Kita menggunakan pendekatan ini untuk meminimalkan bagian-bagian yang menjemukan, yang tidak dapat dihindari dalam landasan yang sepenuhnya aksiomatik.

\newpage
\vspace*{3\baselineskip}

\section{Aksioma-aksioma}
\label{sec:axioms}


\begin{framed}
\begin{enumerate}[I.]
\item\label{def:birkhoff-axioms:0} \index{bidang!bidang Euklides}\index{Euklides!bidang}\emph{Bidang Euklides} adalah ruang metrik dengan sedikitnya dua titik.


\item\label{def:birkhoff-axioms:1} 
Untuk setiap dua titik berbeda $P$ dan $Q$ yang diberikan pada bidang Euklides, terdapat tepat satu garis yang memuat keduanya.

\item\label{def:birkhoff-axioms:2} 
Setiap sudut $\angle AOB$ pada bidang Euklides
menentukan suatu bilangan real dalam interval $(-\pi,\pi]$.
Bilangan ini disebut \index{sudut!ukuran}\emph{ukuran sudut $\angle AOB$}
dan dilambangkan dengan $\measuredangle A O B$.
Bilangan tersebut memenuhi syarat-syarat berikut:
\begin{enumerate}[(a)]
\item\label{def:birkhoff-axioms:2a} 
Jika diberikan suatu sinar $[O A)$ dan $\alpha\in(-\pi,\pi]$, 
terdapat tepat satu sinar $[O B)$ 
sedemikian sehingga $\measuredangle A O B= \alpha$.
\item\label{def:birkhoff-axioms:2b} 
Untuk sebarang titik $A$, $B$, dan $C$ yang berbeda dari $O$, berlaku
$$\measuredangle A O B+\measuredangle B O C
\equiv\measuredangle A O C.$$
\item\label{def:birkhoff-axioms:2c} 
Fungsi 
$$\measuredangle\:(A,O,B)\mapsto\measuredangle A O B$$
kontinu pada sebarang tripel titik $(A,O,B)$
sedemikian sehingga $O\ne A$, $O\ne B$, dan $\measuredangle A O B\ne\pi$.

\end{enumerate}

\item\label{def:birkhoff-axioms:3}  
Pada bidang Euklides, berlaku
$\triangle A B C\cong\triangle A' B' C'$
jika dan hanya jika 
\begin{align*}
A' B'&=A B, & A' C'&= A C, &&\text{dan}
&\measuredangle C' A' B'&=\pm\measuredangle C A B.
\end{align*}
\item\label{def:birkhoff-axioms:4}
Jika untuk dua segitiga $\triangle ABC$, $\triangle AB'C'$ pada bidang Euklides
dan untuk $k>0$ berlaku
\begin{align*}
B'&\in [AB),
& C'&\in [AC),
\\
AB'&=k\cdot AB,&
AC'&=k\cdot AC,
\end{align*}
maka
\begin{align*}
B'C'&=k\cdot BC,&
\measuredangle AB'C'&=\measuredangle ABC,
&
\measuredangle AC'B'&=\measuredangle ACB.
\end{align*}
\end{enumerate}
\end{framed}

Mulai sekarang,  
kita tidak boleh menggunakan informasi apa pun tentang bidang Euklides yang tidak dapat diturunkan dari kelima aksioma di atas.

\begin{thm}{Latihan kelas}\label{ex:infinite}
Buktikan: bidang memuat tak berhingga banyak (a) titik dan (b) garis.
\end{thm}

\section{Garis dan sinar}

\begin{thm}[\abs]{Proposisi}\label{lem:line-line}
\let\thefootnote\relax\footnotetext{${}^\a$ Tanda ``$\a$'' menunjukkan bahwa Aksioma~\ref{def:birkhoff-axioms:4} tidak digunakan dalam bukti.
Abaikan tanda ini untuk saat ini; tanda tersebut akan menjadi penting pada Bab~\ref{chap:non-euclid}.}
Sebarang dua garis berbeda berpotongan pada paling banyak satu titik.
\end{thm}

\parit{Bukti.}
Andaikan dua garis $\ell$ dan $m$ berpotongan di dua titik berbeda $P$ dan~$Q$.
Dengan menerapkan Aksioma~\ref{def:birkhoff-axioms:1}, kita memperoleh $\ell=m$.
\qeds

\begin{thm}{Latihan}\label{ex:[OA)=[OA')}
Misalkan $A'\in[OA)$ dan $A'\not=O$. 
Tunjukkan bahwa 
\[[O A)\z=[O A').\]

\end{thm}

\section{Sudut nol}

\begin{thm}[\abs]{Proposisi}\label{lem:AOA=0}
$\measuredangle A O A= 0$ untuk setiap $A\not=O$.
\end{thm}

\parit{Bukti.}
Menurut Aksioma~\ref{def:birkhoff-axioms:2b},
$$\measuredangle A O A
+
\measuredangle A O A 
\equiv
\measuredangle A O A.$$
Jika $\measuredangle A O A$ dikurangkan dari kedua ruas, diperoleh
$\measuredangle A O A \equiv 0$.

Menurut Aksioma~\ref{def:birkhoff-axioms:2}, $-\pi<\measuredangle A O A\le \pi$;
oleh karena itu, $\measuredangle A O A \z= 0$.
\qeds

\begin{thm}[!]{Latihan}\label{ex:[OA)=[OB)}
Andaikan $\measuredangle A O B= 0$.
Tunjukkan bahwa $[OA)=[OB)$.
\end{thm}

\begin{thm}[\abs]{Proposisi}\label{lem:AOB+BOA=0}
Untuk sebarang titik $A$ dan $B$ yang berbeda dari $O$,
berlaku
$$\measuredangle A O B\equiv-\measuredangle B O A.$$

\end{thm}

\parit{Bukti.}
Menurut Aksioma~\ref{def:birkhoff-axioms:2b},
$$\measuredangle A O B+\measuredangle B O A \equiv\measuredangle A O A.$$
Menurut Proposisi~\ref{lem:AOA=0}, $\measuredangle A O A=0$.
Dengan demikian, hasil tersebut diperoleh.
\qeds

\section{Sudut lurus}

Jika $\measuredangle A O B=\pi$,
kita mengatakan bahwa $\angle A O B$ merupakan 
\index{sudut!sudut lurus}\emph{sudut lurus}.
Menurut Proposisi~\ref{lem:AOB+BOA=0}, 
jika $\angle A O B$ lurus,
maka $\angle B O A$ juga lurus.

Kita mengatakan bahwa titik $O$ \index{di antara}\emph{terletak di antara} titik $A$ dan $B$
jika $O\not= A$, $O\not= B$, dan $O\in[A B]$.

\begin{thm}[\abs]{Teorema}\label{thm:straight-angle}
Sudut $A O B$ merupakan sudut lurus 
jika dan hanya jika $O$ 
terletak di antara $A$ dan~$B$.
\end{thm}

\begin{wrapfigure}{r}{40mm}
\centering
\vskip-0mm
\includegraphics{mppics/pic-8}
\end{wrapfigure}

\parit{Bukti.}
Menurut Latihan~\ref{ex:trig==}, kita dapat mengandaikan bahwa
$O A = O B = 1$.

\parit{Bagian ``jika''.}
Andaikan $O$  
terletak di antara $A$ dan~$B$.
Tetapkan $\alpha=\measuredangle A O B$.

Dengan menerapkan Aksioma~\ref{def:birkhoff-axioms:2a},
kita memperoleh sinar $[OA')$ sedemikian sehingga $\alpha\z=\measuredangle B O A'$.
Menurut Latihan~\ref{ex:trig==}, kita dapat mengandaikan bahwa $OA'=1$.
Menurut Aksioma~\ref{def:birkhoff-axioms:3},
\[\triangle AOB\z\cong\triangle BOA'.\]
Misalkan $f$ menyatakan transformasi isometrik bidang yang bersesuaian;
yaitu, $f$ merupakan transformasi isometrik sedemikian sehingga $f(A)=B$, $f(O)=O$, dan $f(B)=A'$. 

\begin{wrapfigure}{o}{40mm}
\centering
\vskip-2mm
\includegraphics{mppics/pic-10}
\end{wrapfigure}

Maka 
\[(A'B)=f((AB))\ni f(O)=O.\]
Oleh karena itu, kedua garis $(AB)$ dan $(A'B)$ memuat $B$ dan~$O$.
Menurut Aksioma~\ref{def:birkhoff-axioms:1}, $(AB)=(A'B)$.

Menurut definisi garis, $(AB)$ isometrik dengan garis bilangan real $\mathbb{R}$.
Oleh karena itu, $(AB)$ memuat tepat dua titik, $A$ dan $B$, yang berjarak $1$ dari~$O$.
Karena $OA' = 1$, diperoleh $A' = A$ atau $A' = B$.
Karena suatu isometri bersifat injektif, $A' = f(B) \ne f(A) = B$.
Jadi, $A' \ne B$, sehingga $A = A'$.

Menurut Aksioma~\ref{def:birkhoff-axioms:2b} dan Proposisi~\ref{lem:AOA=0}, kita memperoleh
\begin{align*}
2\cdot\alpha&=
\measuredangle AOB+\measuredangle BOA'=
\\
&=\measuredangle AOB+\measuredangle BOA\equiv
\\
&\equiv\measuredangle AOA=
\\
&= 0.
\end{align*}
Oleh karena itu, menurut Latihan~\ref{ex:2a=0}, $\alpha$ sama dengan $0$ atau~$\pi$.

Karena $[OA)\ne [OB)$,  
kita memiliki $\alpha\ne 0$; lihat Latihan~\ref{ex:[OA)=[OB)}.
Oleh karena itu, $\alpha=\pi$.


\parit{Bagian ``hanya jika''.}
Misalkan $\measuredangle A O B= \pi$.
Perhatikan garis $(OA)$ dan pilih titik $B'$ pada $(OA)$ sedemikian sehingga $O$ terletak di antara $A$ dan~$B'$.

Dari hasil di atas, $\measuredangle AOB'=\pi$.
Aksioma~\ref{def:birkhoff-axioms:2a} memberikan $[O B)\z=[O B')$.
Jadi, $O$ terletak di antara $A$ dan~$B$.
\qeds 

Segitiga $ABC$ disebut 
\index{segitiga!segitiga degenerat}\index{degenerat!segitiga}\emph{degenerat}
jika $A$, $B$, dan $C$ terletak pada satu garis.
Akibat berikut hanyalah perumusan ulang Teorema~\ref{thm:straight-angle}.

\begin{thm}[\abs]{Akibat}\label{cor:degenerate=pi}
Suatu segitiga degenerat jika dan hanya jika salah satu sudutnya sama dengan $\pi$ atau~$0$.
Selain itu, ukuran sudut dalam suatu segitiga degenerat adalah $0$, $0$, dan $\pi$.
\end{thm}

\begin{thm}{Latihan}\label{ex:lineAOB}
Tunjukkan bahwa tiga titik berbeda $A$, $O$, dan $B$ terletak pada satu garis jika dan hanya jika 
$$2\cdot \measuredangle AOB\equiv 0.$$ 

\end{thm}

\begin{thm}[!]{Latihan}\label{ex:ABCO-line}
Misalkan $A$, $B$, dan $C$ tiga titik yang berbeda dari~$O$.
Tunjukkan bahwa $B$, $O$, dan $C$ terletak pada satu garis jika dan hanya jika
$$2\cdot \measuredangle AOB\equiv 2\cdot \measuredangle AOC.$$ 

\end{thm}

\begin{thm}{Latihan}\label{ex:infinite-number-of-lines} 
Tunjukkan bahwa terdapat suatu segitiga tak degenerat.
\end{thm}

\section{Sudut bertolak belakang}

Sepasang sudut $AOB$ dan $A'OB'$ 
disebut \index{sudut!sudut bertolak belakang}\index{sudut bertolak belakang}\emph{bertolak belakang}
jika titik $O$ 
terletak di antara $A$ dan $A'$ 
sekaligus di antara $B$ dan $B'$.


\begin{thm}[\abs]{Proposisi}\label{prop:vert}
Sudut-sudut bertolak belakang memiliki ukuran yang sama.
\end{thm}


\parit{Bukti.}
Andaikan sudut $AOB$ dan $A'OB'$ bertolak belakang.
Perhatikan bahwa $\angle AOA'$ dan $\angle BOB'$ lurus.
Oleh karena itu, $\measuredangle AOA'\z=\measuredangle BOB'=\pi$.

{

\begin{wrapfigure}[4]{o}{24mm}
\vskip-2mm
\centering
\includegraphics{mppics/pic-12}
\end{wrapfigure}

Diperoleh
\begin{align*}
0&=\measuredangle AOA'-\measuredangle BOB'\equiv
\\
&\equiv 
\measuredangle AOB+\measuredangle BOA'-\measuredangle BOA'-\measuredangle A'OB'
\equiv
\\
&\equiv\measuredangle AOB-\measuredangle A'OB'.
\end{align*}
Karena $-\pi<\measuredangle AOB\le \pi$ dan $-\pi<\measuredangle A'OB'\le \pi$, kita memperoleh $\measuredangle AOB\z=\measuredangle A'OB'$.
\qeds

}

\section{Pencerminan terhadap suatu titik}

\begin{thm}[\abs]{Lema}\label{lem:point-reflection}
Untuk setiap dua titik berbeda $X$ dan $O$, terdapat tepat satu titik $X'$ sedemikian sehingga $O$ merupakan titik tengah ruas $[XX']$.
\end{thm}

\begin{wrapfigure}{o}{33mm}
\vskip-0mm
\centering
\includegraphics{mppics/pic-915}
\end{wrapfigure}

\parit{Bukti.}
Titik $O$ merupakan titik tengah ruas $[XX']$ jika dan hanya jika $X'\in [XO)$ dan $XX'=2\cdot XO$.

Menurut Aksioma~\ref{def:birkhoff-axioms:1}, garis $(XO)$ dan dengan demikian sinar $[XO)$ ditentukan secara tunggal.
Selanjutnya, menurut Latihan~\ref{ex:trig==}, terdapat tepat satu titik $X'\in [XO)$ sedemikian sehingga $XX'=2\cdot XO$.
\qeds

Titik $X'$ yang diberikan oleh lema disebut \index{pencerminan terhadap suatu titik}\emph{cerminan} titik $X$ terhadap $O$.
Kita juga menetapkan bahwa $O'=O$; yaitu, $O$ merupakan cerminan dirinya sendiri terhadap dirinya sendiri.

Titik $O$ disebut \index{pusat!pencerminan}\emph{pusat pencerminan}.

\begin{thm}[\abs]{Proposisi}\label{prop:point-reflection}
Pencerminan terhadap suatu titik merupakan transformasi isometrik bidang.
\end{thm}

\parit{Bukti.}
Misalkan $O$ pusat pencerminan.
Perhatikan bahwa jika $X'$ merupakan cerminan $X$ terhadap $O$,
maka $X$ merupakan cerminan $X'$.
Dengan kata lain, komposisi pencerminan tersebut dengan dirinya sendiri adalah pemetaan identitas.
Secara khusus, pencerminan tersebut merupakan bijeksi.

\begin{wrapfigure}{o}{33mm}
\vskip-6mm
\centering
\includegraphics{mppics/pic-76}
\end{wrapfigure}

Sekarang pilih dua titik $X$ dan $Y$;
misalkan $X'$ dan $Y'$ secara berurutan merupakan cerminan kedua titik tersebut terhadap $O$.
Untuk memeriksa bahwa pencerminan tersebut mempertahankan jarak, kita perlu menunjukkan bahwa $X'Y'=XY$.

Kita dapat mengandaikan bahwa $X$, $Y$, dan $O$ berbeda; jika tidak, pernyataan tersebut jelas.
Menurut definisi pencerminan, berlaku $OX=OX'$, $OY=OY'$.
Perhatikan pula bahwa sudut $XOY$ dan $X'OY'$ bertolak belakang.
Menurut \ref{prop:vert}, $\measuredangle XOY\z=\measuredangle X'OY'$.

Aksioma~\ref{def:birkhoff-axioms:3} memberikan $\triangle XOY\cong\triangle X'OY'$ dan $X'Y'=XY$.
\qeds

\begin{thm}{Latihan}\label{ex:refelection-of-line}
Buktikan bahwa pencerminan terhadap titik memetakan setiap garis ke garis.
\end{thm}
